\documentclass[conference]{IEEEtran}
\usepackage[utf8]{inputenc}
\usepackage[T1]{fontenc}
\usepackage{cite}
\usepackage{amsmath,amssymb}
\usepackage{graphicx}
\usepackage{booktabs}
\usepackage{url}

\title{Validator‑Guided Repair on a Shared Initial LLM Candidate: A Preregistered Paired Evaluation for Network Configuration Safety}

\author{
\IEEEauthorblockN{Autonomous AI Researcher}
\IEEEauthorblockA{CNSM 2026 Agentic AI Researcher Track}
}

\begin{document}

\maketitle

\begin{abstract}
We preregistered and executed a paired confirmatory experiment (PREREG‑CAND‑001‑VER‑GVR‑2026‑09‑01) to measure whether permitting a deterministic network‑configuration validator plus one automated repair changes the rate of validator‑valid executable configurations relative to leaving the identical sampled LLM candidate unguarded (shared\_initial\_candidate semantics). Using the declared generator gpt‑5‑mini (hosted adapter openai\_responses) and deterministic\_netops\_validator\_v5 on N = 200 paired intents (completed\_episode\_count = 400; model\_calls\_used = 200), every shared initial candidate passed validator checks and no repairs were attempted or applied (repair\_attempt\_count = 0; repair\_applied\_count = 0). Observed outcomes: baseline\_success = 200/200, guarded\_success = 200/200; paired contingency n11 = 200, n10 = 0, n01 = 0, n00 = 0; exact McNemar two‑sided p = 1.0; paired bootstrap percentile 95\% CI (seed = 1729; 10,000 resamples) = [0.0, 0.0]. We present artifact‑grounded methodology, execution accounting, representative validator diagnostics, compact contingency and repair summaries, and a focused discussion clarifying the causal limits of the executed estimand under shared\_initial\_candidate semantics and preregistered follow‑on designs required to measure repair efficacy under independent sampling, higher difficulty, or adversarial inputs.
\end{abstract}

\section{Introduction}
Generative large language models (LLMs) are increasingly proposed to translate operator intent into device configuration sequences in network operations (NetOps). Erroneous sequences, syntactic mistakes, or fabricated fields can cause service disruption or policy violations when applied to devices. A common architectural mitigation is a generate\ensuremath{\rightarrow}validate\ensuremath{\rightarrow}repair (GVR) pipeline: generate candidate configuration(s), deterministically validate them against intent/topology/workflow constraints, and trigger automated repair when the validator finds faults. Prior work documents verifier‑guided improvements in infrastructure‑as‑code and related code‑generation domains and recommends grounding strategies to reduce hallucination in operational tasks [1--3]. However, whether verifier‑guided pipelines reduce execution‑level operational risk for network configurations remains an open empirical question because prior demonstrations often measure textual correctness or IaC test suites rather than executed device‑state safety in packet‑level or emulator contexts [3,4].

This paper reports a fully preregistered paired confirmatory experiment that isolates the downstream causal contribution of a deterministic validator and any permitted repair acting on the exact same sampled LLM candidate (shared\_initial\_candidate semantics). Under these semantics baseline and guarded episodes for each intent begin from the identical initial generation; any paired difference can therefore only arise from deterministic validation and a permitted repair of that same candidate. Our primary research question is: under shared\_initial\_candidate semantics, does permitting a deterministic validator plus at most one automated repair change the proportion of validator‑valid executable configurations relative to leaving the same initial candidate unguarded? The contribution is an artifact‑grounded report of the executed estimand with complete execution accounting, exact inferential statistics, representative validator diagnostics, and precise recommendations for preregistered follow‑on designs that can measure repair efficacy when independent sampling, elevated difficulty, or adversarial inputs are employed.

\section{Related Work}
Hallucination and factual unreliability in LLM outputs are well documented and motivate grounding, retrieval, and verification strategies for operational tasks [1]. Several recent surveys synthesize these failure modes and recommend evidence‑grounding or verification as primary mitigations; they emphasize that hallucination manifests across domains and that mitigation effectiveness depends strongly on task framing and the grounding mechanism used [1].

Retrieval‑augmented and evidence‑grounded pipelines have been shown to reduce hallucination in operations‑style question answering and documentation tasks, primarily by constraining the model with external factual context and by post‑generation grounding checks [2]. These evaluations typically measure textual fidelity (e.g., citation correctness, answer grounding) rather than executable device‑state safety. As a result, reductions in text‑level hallucination do not directly imply reductions in execution risk when generated artifacts are translated into device CLI and applied to network state.

Generate\ensuremath{\rightarrow}validate\ensuremath{\rightarrow}repair (GVR) architectures have been proposed and demonstrated in adjacent domains such as infrastructure‑as‑code (IaC) and program synthesis, where deterministic or policy‑guided verifiers detect rule violations and either supply feedback to the generator or guide repair loops [3]. TerraFormer and related verifier‑guided IaC systems operationalize a feedback loop that fine‑tunes or constrains generation via verifier signals; reported gains are measured against IaC test suites and policy checks rather than execution on heterogeneous device families. Key methodological elements highlighted by these works include (a) verifier expressivity (what constraints and dependency rules are implemented), (b) repair strategy (single corrective edit, multi‑step rewriting, or fine‑tuning), and (c) how verifier feedback is integrated (in‑prompt, tool call, or offline retraining) [3].

NetOps‑specific evaluations of LLM capability concentrate on intent translation accuracy, syntax compliance, and component‑level parsing rather than end‑to‑end execution safety in simulators or testbeds [4]. These studies typically report metrics such as command‑level BLEU/similarities, slot accuracy, or human adjudication of intent fidelity; they provide useful capability baselines but leave open how textual errors map to executable misconfigurations or service disruption when applied to device state. AIOps surveys further underscore the gap between capability benchmarking and execution‑aware evaluation: they call for standardized execution‑level benchmarks and validator fidelity measurements that tie generator outputs to concrete operational safety outcomes [5].

From a methodological standpoint, prior literature therefore separates three evaluation axes: (1) generator sampling variability (how often first‑pass candidates are correct), (2) validator fidelity (whether the verifier captures the operational invariants and dependency rules relevant to execution), and (3) repair efficacy (whether permitted repairs transform invalid candidates into valid, executable ones). Many published demonstrations conflate (1) and (3) by sampling independently across conditions or by evaluating repaired outputs only at the text level. The paired, shared\_initial\_candidate semantics used in our preregistration explicitly isolate axis (2) and (3) acting on a single sampled candidate: baseline and guarded arms begin from the identical generation so any difference must come from deterministic validation and permitted repair rather than from independent sampling.

This distinction clarifies why IaC results are informative but not decisive for NetOps execution safety. TerraFormer‑style verifier‑guided improvements (IaC) indicate that verifier feedback can materially improve generator correctness in structured code domains [3]. However, network device configuration introduces additional complexity: vendor CLI differences, runtime protocol interactions, timing‑sensitive sequences, and multi‑device workflows. Existing NetOps capability studies and AIOps surveys document translation ability and recommend validator uplift but, crucially, do not provide standardized execution‑level outcomes that map generator errors to runnable device states---creating the empirical gap that motivated our execution‑focused preregistered study [4,5].

Finally, the literature highlights practical evaluation considerations that our experiment operationalizes: use deterministic task manifests and contamination checks to avoid corpus leakage, record provider audit traces and model configuration details (including explicitly noting an unset temperature rather than assuming deterministic sampling), and archive per‑episode validator traces (validation\_before/validation\_after, parsed/required command counts, repair\_applied flags) to permit exact replay and independent verification [2--5]. These provenance and scoring practices are commonly recommended but rarely presented end‑to‑end in a preregistered execution bundle that also enforces a narrowly scoped estimand---our study adopts these artifact practices to make the estimand and its limits explicit and reproducible.

\section{Methodology}
Preregistration, estimand, and paired semantics. We preregistered the study as PREREG‑CAND‑001‑VER‑GVR‑2026‑09‑01 and froze the primary estimand before any model calls. Primary estimand: paired\_success\_rate\_difference\_guarded\_minus\_baseline (the per‑intent paired difference in the proportion of validator‑valid executable configurations). The design is a confirmatory paired binary comparison over task\_count = 200 intents producing planned\_episode\_count = 400 episodes. Crucially, generation semantics were registered and executed as shared\_initial\_candidate: exactly one sampled initial generation per pair was produced and deterministically reused by both baseline and guarded episodes. This design choice isolates any downstream causal effect to deterministic validation and the permitted repair acting on the identical candidate; it intentionally removes sampling variability between conditions. We emphasize that the archived model\_configuration.json records temperature = null/unset; null/unset is not evidence of deterministic sampling and does not imply temperature = 0.

Generation and provider audit. The declared generator used in execution was gpt‑5‑mini via the hosted adapter openai\_responses. The execution manifest records model\_calls\_used = 200 (one shared initial generation per pair), hosted\_model\_call\_event\_count = 200, completed\_hosted\_model\_call\_event\_count = 200, failed\_hosted\_model\_call\_event\_count = 0, cache\_miss\_count = 200, and stage\_counts indicating shared\_initial\_generation = 200. These provider‑audit counts confirm that the adapter produced a single sampled candidate per pair for reuse by both arms as preregistered. The experiment permitted no retrieval augmentation (RAG=false in the preregistration) and the adapter enforced the maximum\_model\_calls limit in the preregistration.

Task generation, contamination controls, and task manifest. A deterministic task generator produced the 200 synthetic intent\_configuration\_repair\_v1 tasks. Each manifest entry includes initial\_state, expected\_final\_state, required\_sequence, difficulty annotations (e.g., level tags: very\_hard, extreme), allowed\_touched\_fields, and a generation\_seed to permit exact replay. Preexecution contamination controls enforced an exact‑match SHA‑256 policy and high‑similarity n‑gram checks (n‑gram Jaccard threshold, preregistered high‑similarity flag threshold = 0.9). No exact‑match contamination exclusions were flagged and no confirmatory exclusions were required; any flagged items would have been moved to an exploratory leaked bucket per the contamination\_plan. Contamination artifacts and the fixed task manifest are archived for audit and replay.

Deterministic validator, repair policy, and scored outcome mapping. Scoring used deterministic\_netops\_validator\_v5 with scoring\_rule = "full\_intent\_constraint\_validation." The validator is deterministic and structured: it parses candidate command sequences into normalized\_configuration, compares parsed\_command\_count to task required\_command\_count, enforces workflow\_policies (final\_group\_constraints, minimum\_active\_in\_groups, must\_be\_down\_for\_fields, protected\_interfaces), evaluates dependency rules, and emits structured validator\_trace objects including validation\_before and validation\_after subtables, violation\_codes, violation\_count, parsed\_command\_count, required\_command\_count, observed\_touched\_field\_count, and a repair\_applied boolean. For guarded episodes a single automated repair attempt was permitted (maximum\_repair\_calls\_per\_task = 1); the repair module, if invoked, would generate an edited candidate to be revalidated. Under shared\_initial\_candidate semantics the baseline arm deterministically reused the same sampled candidate without repair; only the guarded arm could alter that identical candidate by invoking the permitted repair. For analysis, validator.valid == true mapped to a binary "executable/valid" outcome per episode; validator.valid == false would trigger repair\_attempt logic in guarded episodes.

Missingness, exclusions, and analysis procedures. The preregistration prescribed the complete\_pair\_primary missingness policy: exclude incomplete pairs from the primary confirmatory McNemar test. A sensitivity analysis (failure‑as‑zero) was preregistered to treat any missing or failed episode as unsafe. The primary hypothesis used an exact two‑sided McNemar test on paired binary outcomes; paired bootstrap percentile CIs were computed (bootstrap\_seed = 1729; bootstrap\_resamples = 10,000) to quantify uncertainty in the paired difference. Exploratory subgroup comparisons by difficulty pattern and failure category were labeled exploratory and would be adjusted using Benjamini--Hochberg FDR = 0.05 when reported. Multiplicity for the main confirmatory test was controlled by design (single prespecified primary test).

Execution reconciliation and reproducibility artifacts. Execution completed without infrastructure failures and matched preregistered accounting: completed\_episode\_count = 400, failed\_episode\_count = 0, complete\_pairs = 200. Deterministic reconciliation checks verified pair accounting, marginal consistency, and provider audit stages. Per‑episode scoring JSON artifacts record the shared\_initial\_candidate, raw\_response\_sha256, normalized\_response, the validator\_trace object, repair\_applied flags, and the scorer\_id/version. Representative artifact excerpts (archived) demonstrate the validator fields used to compute binary outcomes (e.g., parsed\_command\_count, required\_command\_count, normalized\_configuration, violation\_count). All artifacts required to replay the executed estimand (task manifest, responses, scoring traces, analysis CSVs, and model configuration) are archived in the evidence bundle to permit deterministic replay of the exact shared\_initial\_candidate sampling and subsequent scoring.

Implementation notes and limitations of methods. The validator performs static, deterministic checks against task workflow\_policies and expected command sequences; it does not perform vendor CLI execution nor packet‑level emulation. This design choice was deliberate to isolate syntactic and workflow‑policy correctness; it is a recognized construct validity limitation for dynamic runtime behaviors. The preregistered repair policy limited automated repairs to one attempt per guarded episode to bound repair budget and analysis interpretation. Finally, the model\_configuration records temperature = null/unset; the methods explicitly note that null/unset is not evidence of deterministic sampling and that sampling nondeterminism may still have occurred during the single sampled generation that was archived and deterministically reused.

\section{Results}
Execution accounting and deterministic reconciliation. Execution completed per the manifest: planned\_episode\_count = 400, completed\_episode\_count = 400, failed\_episode\_count = 0, model\_calls\_used = 200 (execution/manifests archived). Analysis missingness\_summary reports complete\_pairs = 200 and zero failed or unscorable episodes. Deterministic reconciliation shows stage\_counts = \{shared\_initial\_generation: 200\} and cross\_condition\_cache\_key\_reuse\_observed = false; all deterministic consistency checks passed.

Primary outcomes and contingency. Condition summaries: baseline\_successes = 200/200 and guarded\_successes = 200/200 (success\_rate = 1.0 each). Paired contingency: n11 = 200, n10 = 0, n01 = 0, n00 = 0. Exact McNemar two‑sided p = 1.0. Paired bootstrap (seed = 1729; 10,000 resamples) percentile 95\% CI for the paired difference = [0.0, 0.0]. These artifacts are archived as analysis/condition\_summary.csv and analysis/paired\_contingency\_table.csv.

Repair activity and validator diagnostics. Repair\_attempt\_count = 0 and repair\_applied\_count = 0 across the confirmatory sample: the validator flagged no violations on any shared initial candidate and therefore the permitted guarded repair path was never invoked. Representative per‑episode JSON artifacts (examples archived under execution/scoring/, e.g., task‑000004‑baseline.json and task‑000004‑guarded.json) show identical shared\_initial\_candidate content and identical validation\_after.normalized\_configuration values. Example shared initial candidate (task‑000004):

interface edge2 admin down
interface edge2 vlan 40
interface edge2 admin up
interface edge1 admin down

For that example parsed\_command\_count = 4, required\_command\_count = 4, and validator.valid = true. Similar diagnostics hold for sampled representative tasks (task‑000005, task‑000006).

Contamination and provider audit summary. Preexecution contamination checks flagged no exact matches; analysis/contamination\_summary.csv is empty. Provider audit recorded hosted\_model\_call\_event\_count = 200 and cache\_miss\_count = 200. No infrastructure failures or retrials occurred; analysis/analysis\_log.jsonl documents the bootstrap settings used (bootstrap\_resamples = 10000; bootstrap\_seed = 1729) and that paired analysis completed successfully.

Compact repair summary (explicit). Table 2 (below) provides the preregistered reviewer request summary derived from archived artifacts: total\_pairs = 200; repair\_attempt\_count = 0; repair\_applied\_count = 0. The absence of any validator detections fully explains the observed degenerate contingency: with no validator detections, the guarded repair path had no opportunity to alter outcomes.

\section{Discussion}
Interpreting the executed estimand. The shared\_initial\_candidate semantics were intentionally chosen to isolate the causal effect of deterministic validation and any permitted repair acting on the exact same generator output. The executed estimand therefore answers a narrowly defined causal question: when baseline and guarded episodes begin from the identical sampled candidate, can deterministic validation and at most one automated repair change the validator‑valid/executable outcome? The observed null effect (paired difference = 0.0) shows that, for this task bank and this single sampled candidate per pair, the validator found no violations and the permitted repair path was never exercised. Importantly, this does not evaluate whether guarded pipelines reduce first‑pass generator error rates under independent sampling or different model temperatures; such claims would require a different preregistered estimand and execution.

Artifact‑level diagnostics explaining the degenerate outcome. Multiple archived artifacts explain why the guarded path was not exercised. Representative per‑episode scorer JSONs under execution/scoring/ (for example task‑000004, task‑000005, task‑000006) show that the shared\_initial\_candidate content exactly satisfied the recorded task required\_sequence and workflow\_policies: parsed\_command\_count equaled required\_command\_count in these examples and validation\_after.normalized\_configuration matched the reference expected\_final\_state. The validator traces consistently set validator.valid = true and repair\_applied = false for the sampled examples, and analysis/condition\_summary.csv reports 200 successes in each condition. Provider audit artifacts show hosted\_model\_call\_event\_count = 200 and cache\_miss\_count = 200, confirming that each pair used a single freshly returned initial generation (shared) rather than multiple independent samples. These concrete, archived observations together account for the absence of discordant pairs.

Statistical and inferential implications. From a statistical standpoint the complete absence of discordant pairs (n10 = n01 = 0) makes the confirmatory McNemar test degenerate: there is no observed evidence that the validator or repair step altered any outcome under the executed estimand. This outcome also limits the inferential information the design can provide about repair efficacy; power calculations conducted pre‑execution assumed a nonzero baseline failure rate and thus cannot be used post‑hoc to infer a repair effect here. The artifact record supports this limitation explicitly (analysis/paired\_contingency\_table.csv and analysis/analysis\_log.jsonl). Consequently, meaningful estimation of repair benefit requires a design that induces or selects for nontrivial validator failure rates (see preregistered follow‑on designs below).

Design‑mechanism interpretation without overreach. It is important to avoid two incorrect readings of the result. First, the null paired difference does not imply that validator‑guided repair is unnecessary generally; it only shows that for the executed estimand---one shared sampled candidate per intent drawn from the provided deterministic task bank---the validator detected no faults and therefore repair could not act. Second, the result does not imply that sampling nondeterminism was absent: model\_configuration.json records temperature = null/unset and, as we explicitly note, null/unset is not evidence of deterministic sampling and does not imply temperature = 0. The single shared initial generation per pair was the artifact used in both arms; nondeterministic sampling could still affect a different estimand that samples independently per condition.

Practical lessons for future preregistered experiments. The archived execution suggests concrete, preregistered modifications to exercise the guarded repair path: (1) remove shared\_initial\_candidate semantics so baseline and guarded arms receive independent first‑pass samples (this allows repairs to act on generator faults but requires a revised analysis plan to separate generator variability from repair effects); (2) select or synthesize tasks annotated as "extreme" difficulty or inject adversarial perturbations to raise expected validator failure rates and thereby exercise repair logic (the task manifest includes difficulty annotations that can be used deterministically); (3) preregister non‑null temperature sweeps to increase sample diversity and raise the probability of first‑pass faults; (4) increase validator fidelity by integrating vendor CLI semantics or packet‑level emulation so repairs are evaluated against dynamic runtime behavior rather than only logical workflow constraints; and (5) broaden the model scope beyond a single declared generator to assess whether guard‑repair interactions generalize across generator families. All of these modifications are consistent with the preregistered follow‑on designs archived alongside the execution artifacts and do not require inventing new infrastructure beyond what is already documented.

Validator and task‑bank interactions as a methodological point. The present artifacts show that a deterministic task generator combined with an LLM generator configured under the executed semantics can produce many valid candidates for constrained configuration intents---particularly when the task specification includes a concise required\_sequence and the validator enforces deterministic workflow checks. This observation highlights a methodological interaction: the joint distribution induced by (task generator \ensuremath{\times} sampling semantics \ensuremath{\times} model configuration \ensuremath{\times} validator fidelity) determines whether a guarded repair path will be invoked. Careful preregistered control of each of those factors is therefore necessary to produce a confirmatory test of repair efficacy rather than, as here, a test of downstream validator behavior on already‑valid candidates.

Reproducibility, auditability, and limits of evidence. The execution and analysis artifacts (execution/raw\_results.jsonl, execution/task\_manifest.jsonl, execution/scoring/*.json, analysis/*.csv, model\_configuration.json) provide machine‑verifiable traces that reproduce the executed estimand and demonstrate repair non‑occurrence. Preexecution contamination checks produced no exact‑match exclusions. Deterministic reconciliation and provider audit artifacts show consistent episode accounting and one shared initial generation per pair. Nevertheless, external validity is limited: the task bank is synthetic, the validator is an abstraction that does not execute vendor CLI or perform packet‑level emulation, and the run used a single declared generator and hosted adapter. These constraints are documented in the limitations section and should guide interpretation.

Summary takeaway. The artifact‑grounded evidence is internally consistent and supports a precise, bounded conclusion: under the preregistered shared\_initial\_candidate semantics and the executed synthetic task bank, deterministic validation found no violations and the permitted automated repair path was never exercised, yielding identical valid outcome proportions in baseline and guarded arms. This narrow, reproducible finding clarifies an estimand boundary and motivates preregistered extensions that intentionally increase generator‑level faults or validator fidelity to measure repair utility in operationally challenging regimes.

\section{Limitations}
\begin{itemize}
\item Shared\_initial\_candidate semantics: the design produced exactly one sampled initial generation per intent and reused it across baseline and guarded conditions; this measures validator/repair effects only and cannot detect differences caused by independent per‑condition sampling or prompt engineering.
\item Temperature unset (temperature = null) in the archived model configuration: null/unset is not evidence of deterministic sampling and does not imply temperature = 0; sampling nondeterminism may still have occurred during the single sampled generation.
\item Single model and single hosted provider: results reflect gpt‑5‑mini under the archived adapter; generalization to other models or model families is unsupported by these artifacts.
\item Synthetic task bank and validator abstraction: tasks were synthetic 'intent\_configuration\_repair\_v1' items and the deterministic validator is an abstraction; realistic vendor CLI idiosyncrasies, emergent runtime interactions, and hardware effects were not executed on live devices.
\item No observed validator failures: all initial candidates were validator‑valid, so the experiment could not assess the repair step's effectiveness; the null effect therefore reflects the task bank and sampling semantics rather than evidence that GVR is unnecessary in general.
\end{itemize}

\section{Conclusion}
We executed a fully preregistered paired confirmatory experiment that measures the downstream effect of a deterministic validator plus an allowed single automated repair on the exact same sampled LLM candidate per intent (shared\_initial\_candidate semantics). Using the declared generator gpt‑5‑mini and deterministic\_netops\_validator\_v5 on 200 paired intents, every shared initial candidate passed validation and no repairs were attempted or applied; guarded and baseline outcomes were identical for the executed estimand (paired difference = 0.0; McNemar p = 1.0; paired bootstrap CI = [0.0, 0.0]). This artifact‑grounded null result demonstrates an estimand boundary: under shared\_initial\_candidate semantics and the executed task bank the guarded pipeline could not be exercised and therefore could not change outcomes. It does not imply that GVR pipelines are unnecessary generally. Preregistered follow‑on experiments that (1) remove shared\_initial\_candidate semantics, (2) increase task difficulty or inject adversarial inputs, (3) set explicit sampling parameters, and (4) raise validator fidelity are required to empirically evaluate repair efficacy under realistic sampling variability and adversarial conditions. The archived artifacts and reconciliation files enable exact reproduction of the executed estimand and provide a secure base for precise preregistered extensions.

References
[1] A. Alansari and H. Luqman, "Large Language Models Hallucination: A Comprehensive Survey," arXiv preprint arXiv:2510.06265, 2025.
[2] B. Zhang, H. Rao, and D. Zhao, "Evidence‑Grounded RAG for Cloud‑Native DevOps: Hallucination‑Resistant AIOps Question Answering over Private Operations Documents," J. Autom. Cloud‑Native Syst., 2024.
[3] P. Jana et al., "TerraFormer: Automated Infrastructure‑as‑Code with LLMs Fine‑Tuned via Policy‑Guided Verifier Feedback," Proc. ACM Symp. (2026).
[4] Y. Miao et al., "An Empirical Study of NetOps Capability of Pre‑Trained Large Language Models," arXiv:2309.05557, 2023.
[5] L. Zhang et al., "A Survey of AIOps in the Era of Large Language Models," ACM Comput. Surv., 2025.

One‑line reiteration of the primary estimand limit: the preregistered primary estimand intentionally used shared\_initial\_candidate semantics; evaluating per‑condition independent sampling requires a different preregistration and analysis plan (explicitly recommended above).

\section*{Disclosure Statement}
Disclosure (concise): The human Research Director selected the broad topic family (Generative AI / LLMs for NetOps) before the autonomy lock. After the autonomy lock the autonomous system performed literature review, experiment selection, preregistration (PREREG‑CAND‑001‑VER‑GVR‑2026‑09‑01), code generation, execution, deterministic validation, automated analysis, and manuscript drafting without manual human editing of technical content. Declared model used in execution: gpt‑5‑mini via hosted adapter 'openai\_responses' (execution used model\_calls\_used = 200; archived model\_configuration.json records temperature = null/unset; null/unset is not evidence of deterministic sampling and does not imply temperature = 0). Generation semantics executed: shared\_initial\_candidate --- exactly one sampled initial generation per intent was produced and deterministically reused for both baseline and guarded conditions. Validator: deterministic\_netops\_validator\_v5 scored candidates using 'full\_intent\_constraint\_validation' and a single automated repair iteration was permitted in guarded episodes; in this run the validator flagged no violations and no repairs were applied (repair\_attempt\_count = 0; repair\_applied\_count = 0). Execution accounting: completed\_episode\_count = 400 (200 paired intents) completed without preregistration deviations. Master prompt immutable reference: provenance/master\_prompt.txt --- SHA‑256: 1872df1e\allowbreak{}1805d2d9\allowbreak{}6940456c\allowbreak{}a016bd66\allowbreak{}5d1d5196\allowbreak{}add77f5a\allowbreak{}cdf1582b\allowbreak{}b39b15ba. Pre‑lock infrastructure development and rehearsal occurred but pre‑lock artifacts were not used as empirical evidence in this final run. After the autonomy lock no human manually rewrote technical results, manuscript text, figures, or references.

\begin{thebibliography}{99}
\bibitem{ref1} Aisha Alansari, Hamzah Luqman, "Large Language Models Hallucination: A Comprehensive Survey", 2025, 10.48550/arxiv.2510.06265
\bibitem{ref2} Boning Zhang, Hengning Rao, Derek Zhao , "Evidence-Grounded RAG for Cloud-Native DevOps: Hallucination-Resistant AIOps Question Answering over Private Operations Documents", 2024, 10.69987/jacs.2024.40308
\bibitem{ref3} Prithwish Jana, Sam Davidson, Bhavana Bhasker, Andrey Kan, Anoop Deoras, Laurent Callot, "TerraFormer: Automated Infrastructure-as-Code with LLMs Fine-Tuned via Policy-Guided Verifier Feedback", 2026, 10.1145/3786583.3786898
\bibitem{ref4} Yukai Miao, Yu Bai, Li Chen, Dan Li, Haifeng Sun, Xizheng Wang, Ziqiu Luo, Ren, Yanyu, Dapeng Sun, Xiuting Xu, Qi Zhang, Chao Xiang, Xinchi Li, "An Empirical Study of NetOps Capability of Pre-Trained Large Language Models", 2023, 10.48550/arxiv.2309.05557
\bibitem{ref5} Lingzhe Zhang, Tong Jia, Mengxi Jia, Yifan Wu, Aiwei Liu, Yong Yang, Zhonghai Wu, Xuming Hu, Philip S. Yu, Ying Li, "A Survey of AIOps in the Era of Large Language Models", 2025, 10.1145/3746635
\end{thebibliography}

\end{document}
